\begin{theorem}[Simply connected integer-cell lift with an exterior collar]%
    \label{thm:peps_torus_integer_cell_lift}
    \lean{TNLean.PEPS.exists_integerLift_simplyConnected_exteriorCollar_of_isSimplyConnected}
    \leanok
    \uses{thm:peps_torus_region_planar_lift,
        thm:peps_torus_region_exterior_collar}
    Let $R$ be a finite region of a torus with positive integer
    periods, and assume its actual closed-cell realization $X_R$
    is simply connected. There exists an integer lift $L$ of the
    occupied vertices, preserving each internal unit step. Put
    $A=L(R)$ and let $P_A$ be the union of the closed unit squares
    centered at $A$. Then $P_A$ is simply connected, and every point
    of its exterior collar projects outside $X_R$:
    \begin{align}
        \pi\!\left(\bigl(P_A+[-\tfrac34,\tfrac34]^2\bigr)\setminus P_A\right)
        \cap X_R & =\varnothing.
        \label{eq:peps_integer_lift_collar}
    \end{align}
    Both conclusions hold for the same integer lift. This is an
    auxiliary geometric consequence for the block decomposition in
    \cite[Theorem~6.9]{Schuch2010PEPS}, local source lines 1935--1990.
\end{theorem}

\begin{proof}\leanok
    Simple connectedness gives a continuous lift of the covering
    projection over $X_R$. Its image is homeomorphic to $X_R$ and is
    precisely $P_A$, so $P_A$ is simply connected. Distinct period
    copies of the lifted cells are disjoint. The separation of
    integer centers implies that the enlargement by
    $[-\tfrac34,\tfrac34]^2$ misses every other period copy.
    A point in the exterior collar also misses $P_A$ itself; hence it
    cannot project to $X_R$, proving (\ref{eq:peps_integer_lift_collar}).
\end{proof}
